\documentclass[12pt,a4paper]{report}

% ============================================================
% PACKAGES
% ============================================================
\usepackage[utf8]{inputenc}
\usepackage[T1]{fontenc}
\usepackage[english]{babel}
\usepackage[margin=2.5cm]{geometry}
\usepackage{graphicx}
\usepackage{xcolor}
\usepackage{hyperref}
\usepackage{listings}
\usepackage{fancyhdr}
\usepackage{titlesec}
\usepackage{enumitem}
\usepackage{tcolorbox}
\usepackage{booktabs}
\usepackage{longtable}
\usepackage{tabularx}
\usepackage{fontawesome5}
\usepackage{tikz}
\usepackage{etoolbox}
\usepackage{mdframed}
\usepackage{changepage}

% ============================================================
% COLORS
% ============================================================
\definecolor{accentred}{HTML}{A51C30}
\definecolor{darkblue}{HTML}{1a1a2e}
\definecolor{codeblue}{HTML}{0072B2}
\definecolor{codebg}{HTML}{f5f5f5}
\definecolor{springgreen}{HTML}{6DB33F}
\definecolor{angularred}{HTML}{DD0031}
\definecolor{warningyellow}{HTML}{F5A623}
\definecolor{hintblue}{HTML}{3498DB}
\definecolor{successgreen}{HTML}{27AE60}
\definecolor{lightgray}{HTML}{EEEEEE}
\definecolor{darkgray}{HTML}{333333}
\definecolor{javaorange}{HTML}{F89820}

% ============================================================
% HYPERREF
% ============================================================
\hypersetup{
    colorlinks=true,
    linkcolor=accentred,
    urlcolor=codeblue,
    citecolor=accentred,
    pdftitle={TP - Smart Inventory \& Order Platform},
    pdfauthor={Yassine Elkhamlichi},
    pdfsubject={Full-Stack Angular \& Spring Boot}
}

% ============================================================
% LISTINGS (Code)
% ============================================================
\lstdefinestyle{javastyle}{
    language=Java,
    basicstyle=\ttfamily\small,
    keywordstyle=\color{codeblue}\bfseries,
    commentstyle=\color{gray}\itshape,
    stringstyle=\color{springgreen},
    backgroundcolor=\color{codebg},
    frame=single,
    rulecolor=\color{lightgray},
    numbers=left,
    numberstyle=\tiny\color{gray},
    numbersep=8pt,
    breaklines=true,
    breakatwhitespace=true,
    showstringspaces=false,
    tabsize=4,
    xleftmargin=1.5em,
    framexleftmargin=1em,
    aboveskip=1em,
    belowskip=1em,
}

\lstdefinelanguage{TypeScript}{
    keywords={import, export, from, const, let, var, function, return, if, else, for, while, class, interface, extends, implements, new, this, super, constructor, public, private, protected, static, readonly, async, await, type, enum, true, false, null, undefined, of, in, as, void, number, string, boolean, any, never, unknown, signal, computed, effect, inject},
    sensitive=true,
    comment=[l]{//},
    morecomment=[s]{/*}{*/},
    string=[b]',
    morestring=[b]",
    morestring=[b]`,
}

\lstdefinestyle{tsstyle}{
    language=TypeScript,
    basicstyle=\ttfamily\small,
    keywordstyle=\color{codeblue}\bfseries,
    commentstyle=\color{gray}\itshape,
    stringstyle=\color{springgreen},
    backgroundcolor=\color{codebg},
    frame=single,
    rulecolor=\color{lightgray},
    numbers=left,
    numberstyle=\tiny\color{gray},
    numbersep=8pt,
    breaklines=true,
    breakatwhitespace=true,
    showstringspaces=false,
    tabsize=2,
    xleftmargin=1.5em,
    framexleftmargin=1em,
    aboveskip=1em,
    belowskip=1em,
}

\lstdefinestyle{htmlstyle}{
    language=HTML,
    basicstyle=\ttfamily\small,
    keywordstyle=\color{angularred}\bfseries,
    commentstyle=\color{gray}\itshape,
    stringstyle=\color{springgreen},
    backgroundcolor=\color{codebg},
    frame=single,
    rulecolor=\color{lightgray},
    breaklines=true,
    showstringspaces=false,
    tabsize=2,
    xleftmargin=1.5em,
    framexleftmargin=1em,
    aboveskip=1em,
    belowskip=1em,
}

\lstdefinestyle{yamlstyle}{
    basicstyle=\ttfamily\small,
    keywordstyle=\color{codeblue}\bfseries,
    commentstyle=\color{gray}\itshape,
    stringstyle=\color{springgreen},
    backgroundcolor=\color{codebg},
    frame=single,
    rulecolor=\color{lightgray},
    breaklines=true,
    showstringspaces=false,
    tabsize=2,
    xleftmargin=1.5em,
    framexleftmargin=1em,
    aboveskip=1em,
    belowskip=1em,
}

\lstdefinestyle{dockerstyle}{
    basicstyle=\ttfamily\small,
    keywordstyle=\color{codeblue}\bfseries,
    commentstyle=\color{gray}\itshape,
    stringstyle=\color{springgreen},
    backgroundcolor=\color{codebg},
    frame=single,
    rulecolor=\color{lightgray},
    breaklines=true,
    showstringspaces=false,
    tabsize=2,
    xleftmargin=1.5em,
    framexleftmargin=1em,
    aboveskip=1em,
    belowskip=1em,
}

% ============================================================
% CUSTOM BOXES
% ============================================================
\newtcolorbox{thinkbox}{
    colback=hintblue!5,
    colframe=hintblue,
    title={\faIcon{brain} Think About It},
    fonttitle=\bfseries,
    arc=3mm,
    boxrule=1pt,
    left=8pt,
    right=8pt,
}

\newtcolorbox{checkpoint}{
    colback=successgreen!5,
    colframe=successgreen,
    title={\faIcon{check-circle} Checkpoint},
    fonttitle=\bfseries,
    arc=3mm,
    boxrule=1pt,
    left=8pt,
    right=8pt,
}

\newtcolorbox{hintbox}{
    colback=warningyellow!5,
    colframe=warningyellow,
    title={\faIcon{lightbulb} Hints},
    fonttitle=\bfseries,
    arc=3mm,
    boxrule=1pt,
    left=8pt,
    right=8pt,
}

\newtcolorbox{infobox}{
    colback=codeblue!5,
    colframe=codeblue,
    title={\faIcon{info-circle} Why?},
    fonttitle=\bfseries,
    arc=3mm,
    boxrule=1pt,
    left=8pt,
    right=8pt,
}

% ============================================================
% HEADERS & FOOTERS
% ============================================================
\pagestyle{fancy}
\fancyhf{}
\fancyhead[L]{\small\textcolor{accentred}{\textbf{TP}} --- Smart Inventory Platform}
\fancyhead[R]{\small\textcolor{gray}{Angular \& Spring Boot}}
\fancyfoot[C]{\thepage}
\renewcommand{\headrulewidth}{0.5pt}
\renewcommand{\footrulewidth}{0.3pt}

% ============================================================
% TITLE FORMATTING
% ============================================================
\titleformat{\chapter}[display]
    {\normalfont\huge\bfseries\color{accentred}}
    {\chaptertitlename\ \thechapter}{20pt}{\Huge}
\titleformat{\section}
    {\normalfont\Large\bfseries\color{darkblue}}
    {\thesection}{1em}{}
\titleformat{\subsection}
    {\normalfont\large\bfseries\color{darkblue}}
    {\thesubsection}{1em}{}

% ============================================================
% DOCUMENT
% ============================================================
\begin{document}

% ============================================================
% TITLE PAGE
% ============================================================
\begin{titlepage}
\begin{center}

\vspace*{1cm}

{\Huge\textcolor{accentred}{\textbf{TP}}}\\[0.3cm]
{\large Smart Inventory \& Order Platform}\\[2cm]

\rule{\textwidth}{1.5pt}\\[0.5cm]
{\fontsize{28}{34}\selectfont\textbf{TP --- Smart Inventory}}\\[0.2cm]
{\fontsize{28}{34}\selectfont\textbf{\& Order Platform}}\\[0.5cm]
\rule{\textwidth}{1.5pt}\\[1.5cm]

{\LARGE Full-Stack Development with}\\[0.5cm]
{\Huge\textcolor{springgreen}{\textbf{Spring Boot 3.x}} \& \textcolor{angularred}{\textbf{Angular 18+}}}\\[2cm]

\begin{tabular}{rl}
\textbf{Author:} & Yassine Elkhamlichi \\[0.3cm]
\end{tabular}

\vfill

\begin{tcolorbox}[colback=accentred!5, colframe=accentred, arc=5mm, boxrule=1.5pt, width=0.85\textwidth]
\centering
\textbf{\large Pedagogical Approach --- The 70/30 Golden Rule}\\[0.3cm]
\textbf{30\%} reading documentation \& watching concepts\\
$\downarrow$\\
\textbf{70\%} hands-on coding inside a single repository\\[0.3cm]
\textit{Every concept is immediately applied in the project. No isolated toy examples.}
\end{tcolorbox}

\end{center}
\end{titlepage}

% ============================================================
% TABLE OF CONTENTS
% ============================================================
\tableofcontents
\newpage

% ============================================================
% CHAPTER 1: LEARNING OBJECTIVES
% ============================================================
\chapter{What You Will Learn}

By the end of this TP, you will have built a \textbf{complete, production-grade} full-stack application and mastered the following concepts:

\section{Spring Boot 3.x (Backend)}

\begin{longtable}{p{4.5cm}p{9cm}}
\toprule
\textbf{Concept} & \textbf{Description} \\
\midrule
\endhead
Layered Architecture & Controller $\to$ Service $\to$ Repository --- strict separation of concerns \\[0.3cm]
JPA / Hibernate & Entity mapping, relationships, \texttt{@EntityGraph}, DTO Projections \\[0.3cm]
REST API Design & Versioned endpoints, proper HTTP verbs, status codes \\[0.3cm]
Bean Validation & \texttt{@Valid}, \texttt{@NotBlank}, \texttt{@Min}, custom constraints \\[0.3cm]
Global Error Handling & \texttt{@RestControllerAdvice}, \texttt{@ExceptionHandler}, structured error maps \\[0.3cm]
Pagination \& Sorting & \texttt{Pageable}, \texttt{Page<T>}, dynamic search with JPA Specifications \\[0.3cm]
Spring Security + JWT & Stateless authentication, \texttt{JwtAuthenticationFilter}, role-based access \\[0.3cm]
Performance & N+1 problem diagnosis, \texttt{@EntityGraph}, DTO Projections \\
\bottomrule
\end{longtable}

\section{Angular 18+ (Frontend)}

\begin{longtable}{p{4.5cm}p{9cm}}
\toprule
\textbf{Concept} & \textbf{Description} \\
\midrule
\endhead
Standalone Architecture & No \texttt{NgModule} --- \texttt{provideHttpClient()}, \texttt{provideRouter()} \\[0.3cm]
Strongly-Typed Services & \texttt{HttpClient} with TypeScript interfaces \\[0.3cm]
Modern Control Flow & \texttt{@for}, \texttt{@if}, \texttt{@switch}, \texttt{@defer} \\[0.3cm]
Reactive Forms & \texttt{FormGroup}, \texttt{FormControl}, \texttt{Validators}, custom validators \\[0.3cm]
Angular Signals & \texttt{signal()}, \texttt{computed()}, \texttt{effect()} for reactive state \\[0.3cm]
RxJS Operators & \texttt{debounceTime}, \texttt{distinctUntilChanged}, \texttt{switchMap} \\[0.3cm]
HTTP Interceptors & \texttt{HttpInterceptorFn} for token injection and error handling \\[0.3cm]
Route Guards & \texttt{CanActivateFn} for protecting routes \\[0.3cm]
Lazy Loading & \texttt{@defer (on viewport)} for performance optimization \\
\bottomrule
\end{longtable}

\section{Architecture \& DevOps}

\begin{longtable}{p{4.5cm}p{9cm}}
\toprule
\textbf{Concept} & \textbf{Description} \\
\midrule
\endhead
Clean Project Structure & Industry-standard folder naming used by top tech companies \\[0.3cm]
Full-Stack Integration & Frontend $\leftrightarrow$ Backend communication with proper CORS \\[0.3cm]
Docker Compose & Multi-container deployment (DB + API + UI) \\[0.3cm]
Debugging Skills & Network Tab, SQL logs, error tracing \\
\bottomrule
\end{longtable}

% ============================================================
% CHAPTER 2: PROJECT OVERVIEW
% ============================================================
\chapter{Project Overview \& Architecture}

\section{Smart Inventory \& Order Platform}

You are building a \textbf{professional inventory management system} with three main modules:

\begin{itemize}[leftmargin=2cm]
    \item[\faIcon{box}] \textbf{Products Module:} Live search, category filtering, pagination, sorting
    \item[\faIcon{shopping-cart}] \textbf{Cart Module:} Add/remove items, update quantities, server sync
    \item[\faIcon{lock}] \textbf{Admin Panel:} Role-based access, price management, stock control
\end{itemize}

\section{Backend Architecture (Spring Boot)}

\begin{lstlisting}[style=yamlstyle, title={Spring Boot Project Structure}]
smart-inventory-api/
|-- src/main/java/com/smartinventory/
|   |-- SmartInventoryApplication.java
|   |-- config/
|   |   |-- CorsConfig.java
|   |   +-- SecurityConfig.java
|   |-- controller/
|   |   |-- AuthController.java
|   |   |-- ProductController.java
|   |   |-- CategoryController.java
|   |   |-- CartController.java
|   |   +-- AdminController.java
|   |-- dto/
|   |   |-- request/
|   |   |   |-- LoginRequest.java
|   |   |   |-- RegisterRequest.java
|   |   |   |-- ProductRequest.java
|   |   |   +-- CartItemRequest.java
|   |   +-- response/
|   |       |-- AuthResponse.java
|   |       |-- ProductResponse.java
|   |       |-- PageResponse.java
|   |       +-- ErrorResponse.java
|   |-- exception/
|   |   |-- GlobalExceptionHandler.java
|   |   |-- ResourceNotFoundException.java
|   |   +-- BusinessException.java
|   |-- model/
|   |   |-- Product.java
|   |   |-- Category.java
|   |   |-- CartItem.java
|   |   |-- Order.java
|   |   +-- User.java
|   |-- repository/
|   |   |-- ProductRepository.java
|   |   |-- CategoryRepository.java
|   |   |-- CartItemRepository.java
|   |   +-- UserRepository.java
|   |-- security/
|   |   |-- JwtService.java
|   |   |-- JwtAuthenticationFilter.java
|   |   +-- UserDetailsServiceImpl.java
|   +-- service/
|       |-- ProductService.java
|       |-- CategoryService.java
|       |-- CartService.java
|       |-- AuthService.java
|       +-- impl/
|           |-- ProductServiceImpl.java
|           |-- CategoryServiceImpl.java
|           |-- CartServiceImpl.java
|           +-- AuthServiceImpl.java
+-- pom.xml
\end{lstlisting}

\section{Frontend Architecture (Angular)}

\begin{lstlisting}[style=yamlstyle, title={Angular Project Structure}]
smart-inventory-ui/
+-- src/app/
    |-- core/
    |   |-- interceptors/
    |   |   |-- auth.interceptor.ts
    |   |   +-- error.interceptor.ts
    |   |-- guards/
    |   |   |-- auth.guard.ts
    |   |   +-- admin.guard.ts
    |   +-- services/
    |       |-- auth.service.ts
    |       +-- storage.service.ts
    |-- features/
    |   |-- auth/
    |   |   |-- login/
    |   |   |-- register/
    |   |   |-- models/
    |   |   +-- auth.routes.ts
    |   |-- products/
    |   |   |-- product-list/
    |   |   |-- product-form/
    |   |   |-- models/
    |   |   |-- services/
    |   |   +-- products.routes.ts
    |   |-- cart/
    |   +-- admin/
    |-- shared/
    |   |-- components/
    |   |   |-- navbar/
    |   |   +-- pagination/
    |   +-- pipes/
    |-- app.component.ts
    |-- app.config.ts
    +-- app.routes.ts
\end{lstlisting}

\begin{infobox}
\textbf{Why this structure?} This is the standard used at companies like Google, Apple, and Microsoft.
\begin{itemize}
    \item \texttt{core/} --- Singleton services, guards, interceptors (instantiated ONCE)
    \item \texttt{features/} --- Domain modules, each self-contained
    \item \texttt{shared/} --- Reusable components, pipes, directives
\end{itemize}
\end{infobox}

% ============================================================
% CHAPTER 3: PHASE 1
% ============================================================
\chapter{Phase 1 --- Data Contract \& Basic Rendering}

\textbf{Objective:} Establish the full-stack connection. By the end of this phase, your Angular app will fetch and display a list of products from your Spring Boot API.

\textbf{Estimated Time:} 4 days (2h/day)

% --- Exercise 1.1 ---
\section{Exercise 1.1 --- Initialize the Spring Boot Project}

\textbf{Context:} Every professional project starts with a well-configured foundation.

\textbf{Tasks:}
\begin{enumerate}
    \item Go to \href{https://start.spring.io}{start.spring.io} and generate a project with:
    \begin{itemize}
        \item Group: \texttt{com.smartinventory}, Artifact: \texttt{smart-inventory-api}
        \item Java 21
        \item Dependencies: Spring Web, Spring Data JPA, PostgreSQL Driver, Lombok, Validation
    \end{itemize}
    \item Configure your database connection in \texttt{application.yml}. Think about:
    \begin{itemize}
        \item What is the JDBC URL format for PostgreSQL?
        \item What property controls table auto-creation in JPA?
        \item How do you enable SQL logging?
    \end{itemize}
    \item Create ALL the package folders shown in the Architecture Blueprint.
\end{enumerate}

\begin{thinkbox}
Why do we separate \texttt{dto/request} and \texttt{dto/response} into sub-packages? What problem does this solve compared to putting all DTOs in a single package?
\end{thinkbox}

\begin{checkpoint}
Your Spring Boot application starts without errors, connects to PostgreSQL, and the console shows the Hibernate dialect being used.
\end{checkpoint}

% --- Exercise 1.2 ---
\section{Exercise 1.2 --- Design the Data Model}

\textbf{Context:} The \texttt{model/} package contains your JPA entities. Start with \texttt{Product} and \texttt{Category}.

\textbf{Tasks:}
\begin{enumerate}
    \item Create the \texttt{Category} entity: \texttt{id} (Long, auto-generated), \texttt{name} (String, unique), \texttt{description} (String), and a one-to-many relationship with products.
    \item Create the \texttt{Product} entity: \texttt{id}, \texttt{name}, \texttt{description}, \texttt{price} (BigDecimal), \texttt{stockQuantity} (Integer), \texttt{imageUrl}, \texttt{createdAt}, \texttt{updatedAt}, and a many-to-one relationship with category.
\end{enumerate}

\begin{thinkbox}
\begin{itemize}
    \item Why do we use \texttt{BigDecimal} for price instead of \texttt{Double}?
    \item What JPA annotation makes \texttt{createdAt} auto-populate on first save?
    \item What is the difference between \texttt{FetchType.LAZY} and \texttt{FetchType.EAGER}?
\end{itemize}
\end{thinkbox}

\begin{hintbox}
Look at \texttt{@Entity}, \texttt{@Table}, \texttt{@Id}, \texttt{@GeneratedValue}, \texttt{@CreationTimestamp}, \texttt{@UpdateTimestamp}, \texttt{@ManyToOne}, \texttt{@OneToMany(mappedBy = "...")}
\end{hintbox}

\begin{checkpoint}
Restart the app. Hibernate should auto-create \texttt{products} and \texttt{categories} tables.
\end{checkpoint}

% --- Exercise 1.3 ---
\section{Exercise 1.3 --- Build the Repository Layer}

\textbf{Context:} In Spring Data JPA, you don't write SQL --- you extend an interface.

\textbf{Tasks:}
\begin{enumerate}
    \item Create \texttt{CategoryRepository} extending the appropriate Spring Data interface.
    \item Create \texttt{ProductRepository} with a custom derived query method to find products by category ID.
\end{enumerate}

\begin{thinkbox}
Spring Data JPA uses ``Derived Query Methods.'' If your entity has a field \texttt{category} with a sub-field \texttt{id}, what method name would automatically generate \texttt{WHERE category\_id = ?}?
\end{thinkbox}

% --- Exercise 1.4 ---
\section{Exercise 1.4 --- Create the DTO Layer}

\textbf{Context:} \textbf{Never expose your entities directly in API responses.} DTOs act as a shield between your database model and the outside world.

\textbf{Tasks:}
\begin{enumerate}
    \item Create \texttt{ProductResponse} record with: \texttt{id}, \texttt{name}, \texttt{price}, \texttt{stockQuantity}, \texttt{imageUrl}, \texttt{categoryName}
    \item Create \texttt{CategoryResponse} record with \texttt{id} and \texttt{name}
    \item Create \texttt{ProductRequest} record with all creation fields including \texttt{categoryId}
\end{enumerate}

\begin{thinkbox}
\begin{itemize}
    \item Why do we use Java \texttt{record} instead of a regular class for DTOs?
    \item Why does \texttt{ProductResponse} contain \texttt{categoryName} instead of \texttt{categoryId}?
    \item What would happen if we returned the entity directly with a lazily loaded field?
\end{itemize}
\end{thinkbox}

% --- Exercise 1.5 ---
\section{Exercise 1.5 --- Implement the Service Layer}

\textbf{Context:} The Service layer contains business logic. It receives DTOs from controllers and returns DTOs back.

\textbf{Tasks:}
\begin{enumerate}
    \item Create a \texttt{ProductService} \textbf{interface} with: \texttt{getAllProducts()}, \texttt{getProductById(Long id)}, \texttt{createProduct(ProductRequest request)}
    \item Create \texttt{ProductServiceImpl} implementing the interface:
    \begin{itemize}
        \item Inject repositories via constructor injection
        \item Convert entities to DTOs in each method
        \item Throw \texttt{ResourceNotFoundException} when not found
    \end{itemize}
\end{enumerate}

\begin{thinkbox}
\begin{itemize}
    \item Why do we define an interface separate from the implementation?
    \item Why constructor injection instead of \texttt{@Autowired} on fields?
    \item Should entity-to-DTO conversion live in the service or a separate mapper?
\end{itemize}
\end{thinkbox}

% --- Exercise 1.6 ---
\section{Exercise 1.6 --- Build the REST Controller}

\textbf{Context:} Controllers are the entry point of your API. They should be \textbf{thin} --- no business logic.

\textbf{Tasks:}
\begin{enumerate}
    \item Create \texttt{ProductController} mapped to \texttt{/api/v1/products}
    \item Implement: \texttt{GET /} (list), \texttt{GET /\{id\}} (detail), \texttt{POST /} (create with HTTP 201)
\end{enumerate}

\begin{thinkbox}
\begin{itemize}
    \item Why do we version the API with \texttt{/api/v1/}?
    \item Should the controller return \texttt{ResponseEntity<T>} or just \texttt{T}?
\end{itemize}
\end{thinkbox}

\begin{checkpoint}
Test with Postman: create a product via POST, retrieve all via GET.
\end{checkpoint}

% --- Exercise 1.7 ---
\section{Exercise 1.7 --- Configure CORS}

\textbf{Context:} Your Angular app runs on port 4200, your API on 8080. Browsers block cross-origin requests by default.

\textbf{Tasks:}
\begin{enumerate}
    \item Create \texttt{CorsConfig} implementing \texttt{WebMvcConfigurer}
    \item Allow requests from \texttt{http://localhost:4200}
    \item Allow methods: GET, POST, PUT, DELETE, OPTIONS
\end{enumerate}

\begin{thinkbox}
Why is CORS a server-side configuration? In production, should you allow \texttt{*} (all origins)?
\end{thinkbox}

% --- Exercise 1.8 ---
\section{Exercise 1.8 --- Initialize the Angular Project}

\textbf{Tasks:}
\begin{enumerate}
    \item Generate: \texttt{ng new smart-inventory-ui --standalone --style=css --routing --ssr=false}
    \item Create ALL folders from the Angular Architecture Blueprint
    \item Configure environments with \texttt{apiUrl}
    \item Set up \texttt{app.config.ts} with \texttt{provideHttpClient(withFetch())} and \texttt{provideRouter(routes)}
\end{enumerate}

\begin{thinkbox}
What is the difference between the old NgModule architecture and Standalone? Why \texttt{withFetch()}?
\end{thinkbox}

% --- Exercise 1.9 ---
\section{Exercise 1.9 --- Create TypeScript Models}

\textbf{Tasks:}
\begin{enumerate}
    \item In \texttt{features/products/models/product.model.ts}, create \texttt{Product} and \texttt{ProductRequest} interfaces
    \item Ensure they mirror the backend DTOs exactly
\end{enumerate}

\begin{thinkbox}
Why do we use \texttt{interface} and not \texttt{class}? What happens if the backend adds a field you don't have?
\end{thinkbox}

% --- Exercise 1.10 ---
\section{Exercise 1.10 --- Build the Product Service}

\textbf{Tasks:}
\begin{enumerate}
    \item Create \texttt{product.service.ts} with \texttt{@Injectable(\{ providedIn: 'root' \})}
    \item Implement strongly-typed methods: \texttt{getAll()}, \texttt{getById(id)}, \texttt{create(request)}
\end{enumerate}

\begin{thinkbox}
Why does \texttt{HttpClient.get<Product[]>()} return an \texttt{Observable} instead of a \texttt{Promise}?
\end{thinkbox}

% --- Exercise 1.11 ---
\section{Exercise 1.11 --- Create the Product List Component}

\textbf{Tasks:}
\begin{enumerate}
    \item Generate: \texttt{ng generate component features/products/product-list --standalone}
    \item Inject \texttt{ProductService}, fetch data in \texttt{ngOnInit()}
    \item Use \texttt{@for} (not \texttt{*ngFor}) with \texttt{track} to display products
    \item Use \texttt{@if} to show ``Out of Stock'' conditionally
    \item Use \texttt{@empty} block for empty lists
    \item Set up routing in \texttt{app.routes.ts}
\end{enumerate}

\begin{thinkbox}
What is the \texttt{track} expression in \texttt{@for}? What is the memory leak risk with \texttt{subscribe()} in \texttt{ngOnInit()}?
\end{thinkbox}

\begin{checkpoint}
Navigate to \texttt{http://localhost:4200/products}. Products from your API should appear.
\end{checkpoint}

% --- Exercise 1.12 ---
\section{Exercise 1.12 --- Navigation \& Layout}

\textbf{Tasks:}
\begin{enumerate}
    \item Create \texttt{NavbarComponent} in \texttt{shared/components/navbar/}
    \item Use \texttt{routerLink} (not \texttt{href}) and \texttt{routerLinkActive}
    \item Update \texttt{app.component.html} with navbar and \texttt{<router-outlet>}
\end{enumerate}

\begin{thinkbox}
Why \texttt{routerLink} instead of \texttt{href}? What happens to the page state with \texttt{href}?
\end{thinkbox}

% ============================================================
% CHAPTER 4: PHASE 2
% ============================================================
\chapter{Phase 2 --- Reactive Forms \& Error Handling}

\textbf{Objective:} Build a product creation form with client-side and server-side validation.

% --- Exercise 2.1 ---
\section{Exercise 2.1 --- Backend Validation}

\textbf{Tasks:}
\begin{enumerate}
    \item Add validation annotations to \texttt{ProductRequest}: \texttt{@NotBlank}, \texttt{@Size}, \texttt{@DecimalMin}, \texttt{@Min}, \texttt{@NotNull}
    \item Activate validation with \texttt{@Valid} on the controller parameter
\end{enumerate}

\begin{thinkbox}
What is the difference between \texttt{@NotNull}, \texttt{@NotEmpty}, and \texttt{@NotBlank}?
\end{thinkbox}

% --- Exercise 2.2 ---
\section{Exercise 2.2 --- Global Error Handler}

\textbf{Tasks:}
\begin{enumerate}
    \item Create \texttt{ErrorResponse} record: \texttt{timestamp}, \texttt{status}, \texttt{error}, \texttt{errors} (Map)
    \item Create \texttt{GlobalExceptionHandler} with \texttt{@RestControllerAdvice}:
    \begin{itemize}
        \item Handle \texttt{MethodArgumentNotValidException} $\to$ HTTP 400 with field errors
        \item Handle \texttt{ResourceNotFoundException} $\to$ HTTP 404
        \item Handle generic \texttt{Exception} $\to$ HTTP 500 (never expose stack traces!)
    \end{itemize}
\end{enumerate}

\begin{thinkbox}
Why is \texttt{@RestControllerAdvice} better than try-catch in every controller? What's the difference between what you \emph{log} and what you \emph{return}?
\end{thinkbox}

% --- Exercise 2.3 ---
\section{Exercise 2.3 --- Product Form (Angular)}

\textbf{Tasks:}
\begin{enumerate}
    \item Build a \texttt{FormGroup} with \texttt{FormControl} and \texttt{Validators} for each field
    \item On submit: \texttt{markAllAsTouched()}, validate, then call the service
    \item Display validation errors below each field (only when touched AND invalid)
\end{enumerate}

\begin{thinkbox}
What is ``touched'' vs ``dirty''? Why \texttt{markAllAsTouched()} on submit?
\end{thinkbox}

% --- Exercise 2.4 ---
\section{Exercise 2.4 --- Display Server-Side Errors}

\textbf{Tasks:}
\begin{enumerate}
    \item Create a \texttt{serverErrors} object to store error map from HTTP 400 responses
    \item Display server errors under the corresponding fields
    \item Clear server errors when the user starts typing (listen to \texttt{valueChanges})
\end{enumerate}

% --- Exercise 2.5 ---
\section{Exercise 2.5 --- Custom Validator}

\textbf{Tasks:}
\begin{enumerate}
    \item Create a custom validator function for price format (max 2 decimal places)
    \item Return \texttt{\{ invalidPrice: true \}} if invalid, \texttt{null} if valid
    \item Apply to the price \texttt{FormControl} and display the error message
\end{enumerate}

\begin{thinkbox}
What is the difference between a synchronous and async validator? When would you need async?
\end{thinkbox}

% ============================================================
% CHAPTER 5: PHASE 3
% ============================================================
\chapter{Phase 3 --- Signals \& RxJS Pipeline}

\textbf{Objective:} Implement live search, pagination, and sorting. Master Angular Signals and RxJS operators.

% --- Exercise 3.1 ---
\section{Exercise 3.1 --- Backend Pagination \& Sorting}

\textbf{Tasks:}
\begin{enumerate}
    \item Modify service to accept \texttt{Pageable} and return \texttt{Page<ProductResponse>}
    \item Create \texttt{PageResponse<T>} generic DTO
    \item Update controller to accept \texttt{page}, \texttt{size}, \texttt{sort} query parameters
\end{enumerate}

\begin{thinkbox}
What SQL query does \texttt{Pageable} generate? What is the difference between \texttt{Slice<T>} and \texttt{Page<T>}?
\end{thinkbox}

% --- Exercise 3.2 ---
\section{Exercise 3.2 --- Dynamic Search with JPA Specifications}

\textbf{Tasks:}
\begin{enumerate}
    \item Extend \texttt{ProductRepository} with \texttt{JpaSpecificationExecutor<Product>}
    \item Create static specification methods: \texttt{hasName()}, \texttt{hasCategory()}, \texttt{hasPriceRange()}
    \item Combine specifications with \texttt{.and()} in the service layer
\end{enumerate}

\begin{thinkbox}
How does the Specification pattern implement the Strategy design pattern? What if all parameters are null?
\end{thinkbox}

% --- Exercise 3.3 ---
\section{Exercise 3.3 --- Angular Signals}

\textbf{Tasks:}
\begin{enumerate}
    \item Replace plain variables with \texttt{signal()}: \texttt{products}, \texttt{loading}, \texttt{currentPage}
    \item Create \texttt{computed()} signals for derived state
    \item Update templates to use \texttt{()} syntax: \texttt{products()}, \texttt{loading()}
    \item Use \texttt{.set()} and \texttt{.update()} appropriately
\end{enumerate}

\begin{thinkbox}
What is the difference between \texttt{signal.set()} and \texttt{signal.update()}? How do Signals compare to React's \texttt{useState()}?
\end{thinkbox}

% --- Exercise 3.4 ---
\section{Exercise 3.4 --- Live Search with RxJS}

\textbf{Tasks:}
\begin{enumerate}
    \item Add a search \texttt{FormControl} bound to an input
    \item Build the RxJS pipeline: \texttt{valueChanges} $\to$ \texttt{debounceTime(300)} $\to$ \texttt{distinctUntilChanged()} $\to$ \texttt{switchMap()} $\to$ \texttt{subscribe()}
    \item Clean up with \texttt{takeUntil(destroy\$)} in \texttt{ngOnDestroy()}
\end{enumerate}

\begin{thinkbox}
Why is \texttt{switchMap} the correct operator, not \texttt{mergeMap} or \texttt{concatMap}? What does ``cancel the previous request'' mean at the HTTP level?
\end{thinkbox}

\begin{checkpoint}
Type in the search box. Check the Network tab --- previous requests should be cancelled.
\end{checkpoint}

% --- Exercise 3.5 ---
\section{Exercise 3.5 --- Reusable Pagination Component}

\textbf{Tasks:}
\begin{enumerate}
    \item Create \texttt{PaginationComponent} in \texttt{shared/components/pagination/}
    \item Define inputs: \texttt{currentPage}, \texttt{totalPages}
    \item Define output: \texttt{pageChange} event
    \item Display Previous/Next buttons with page numbers
\end{enumerate}

\begin{thinkbox}
What is a ``smart'' vs ``dumb/presentational'' component? Which is PaginationComponent?
\end{thinkbox}

% ============================================================
% CHAPTER 6: PHASE 4
% ============================================================
\chapter{Phase 4 --- JWT Security \& Interceptors}

\textbf{Objective:} Implement complete authentication and authorization.

% --- Exercise 4.1 ---
\section{Exercise 4.1 --- User Entity \& Auth DTOs}

\textbf{Tasks:}
\begin{enumerate}
    \item Create \texttt{User} entity with: \texttt{id}, \texttt{firstName}, \texttt{lastName}, \texttt{email}, \texttt{password} (hashed), \texttt{role} (enum: USER, ADMIN)
    \item Create \texttt{LoginRequest}, \texttt{RegisterRequest}, \texttt{AuthResponse} DTOs
    \item Create \texttt{UserRepository} with \texttt{findByEmail()} and \texttt{existsByEmail()}
\end{enumerate}

\begin{thinkbox}
Why hashed passwords? Should \texttt{Role} be an entity or enum? Where should password-match validation happen?
\end{thinkbox}

% --- Exercise 4.2 ---
\section{Exercise 4.2 --- JWT Service}

\textbf{Tasks:}
\begin{enumerate}
    \item Add jjwt dependencies to \texttt{pom.xml}
    \item Implement \texttt{JwtService}: \texttt{generateToken()}, \texttt{extractUsername()}, \texttt{isTokenValid()}
    \item Store secret key in \texttt{application.yml}
\end{enumerate}

\begin{thinkbox}
What are the three parts of a JWT? Why is the secret key critical? Symmetric vs asymmetric signing?
\end{thinkbox}

% --- Exercise 4.3 ---
\section{Exercise 4.3 --- Spring Security Configuration}

\textbf{Tasks:}
\begin{enumerate}
    \item Create \texttt{UserDetailsServiceImpl} implementing \texttt{UserDetailsService}
    \item Create \texttt{JwtAuthenticationFilter} extending \texttt{OncePerRequestFilter}
    \item Create \texttt{SecurityConfig}: disable CSRF, set STATELESS sessions, configure URL authorization, add JWT filter
\end{enumerate}

\begin{thinkbox}
Why disable CSRF for a REST API? What does ``stateless'' mean? Why must the JWT filter run before \texttt{UsernamePasswordAuthenticationFilter}?
\end{thinkbox}

% --- Exercise 4.4 ---
\section{Exercise 4.4 --- Auth Service \& Controller}

\textbf{Tasks:}
\begin{enumerate}
    \item Implement \texttt{AuthService}: register (check duplicates, hash password, generate token), login (authenticate, generate token)
    \item Create \texttt{AuthController}: \texttt{POST /auth/register}, \texttt{POST /auth/login}
\end{enumerate}

\begin{checkpoint}
Register $\to$ get token. Login $\to$ get token. Use token to access \texttt{GET /products}.
\end{checkpoint}

% --- Exercise 4.5-4.6 ---
\section{Exercise 4.5 --- Angular Auth \& Interceptors}

\textbf{Tasks:}
\begin{enumerate}
    \item Create \texttt{StorageService}: \texttt{setToken()}, \texttt{getToken()}, \texttt{removeToken()}, \texttt{isLoggedIn()}
    \item Create \texttt{AuthService}: \texttt{login()}, \texttt{register()}, \texttt{logout()}, \texttt{getUserRole()}
    \item Create \texttt{authInterceptor}: clone requests and add \texttt{Authorization: Bearer <token>}
    \item Create \texttt{errorInterceptor}: redirect on 401, handle 403
    \item Register interceptors in \texttt{app.config.ts}
\end{enumerate}

\begin{thinkbox}
Why clone the request? Why skip auth endpoints? Why is \texttt{localStorage} a security risk (XSS)?
\end{thinkbox}

% --- Exercise 4.7-4.8 ---
\section{Exercise 4.7 --- Route Guards \& Auth Pages}

\textbf{Tasks:}
\begin{enumerate}
    \item Create \texttt{authGuard} (\texttt{CanActivateFn}): check login, redirect to \texttt{/auth/login}
    \item Create \texttt{adminGuard}: check ADMIN role
    \item Apply guards to routes
    \item Build \texttt{LoginComponent} and \texttt{RegisterComponent} with Reactive Forms
    \item Implement cross-field validator for password confirmation
\end{enumerate}

\begin{checkpoint}
Complete auth flow: Register $\to$ Login $\to$ Access protected pages $\to$ Logout $\to$ Redirect.
\end{checkpoint}

% ============================================================
% CHAPTER 7: PHASE 5
% ============================================================
\chapter{Phase 5 --- Production Readiness}

\textbf{Objective:} Optimize performance, solve N+1, add lazy loading, containerize with Docker.

\section{Exercise 5.1 --- N+1 Problem}

\textbf{Tasks:}
\begin{enumerate}
    \item Enable SQL logging, count queries for a product list
    \item Fix with \texttt{@EntityGraph(attributePaths = \{"category"\})}
    \item Alternative: DTO Projection with JPQL \texttt{SELECT new}
\end{enumerate}

\begin{thinkbox}
Why is \texttt{FetchType.EAGER} not a good fix? When is DTO Projection better than \texttt{@EntityGraph}?
\end{thinkbox}

\section{Exercise 5.2 --- Angular Deferred Loading}

\textbf{Tasks:}
\begin{enumerate}
    \item Wrap heavy components in \texttt{@defer (on viewport)} blocks
    \item Create skeleton loader CSS animations
    \item Experiment with \texttt{(on idle)}, \texttt{(on timer)}, \texttt{(when condition)}
\end{enumerate}

\section{Exercise 5.3 --- Docker Compose}

\textbf{Tasks:}
\begin{enumerate}
    \item Create multi-stage \texttt{Dockerfile} for Spring Boot (Maven build $\to$ JRE run)
    \item Create multi-stage \texttt{Dockerfile} for Angular (Node build $\to$ Nginx serve)
    \item Create \texttt{nginx.conf} with API proxy and SPA routing
    \item Create \texttt{docker-compose.yml} with 3 services: db, api, ui
\end{enumerate}

\begin{thinkbox}
Why multi-stage builds? Why does Nginx proxy API requests? What is a Docker volume?
\end{thinkbox}

\begin{checkpoint}
\texttt{docker-compose up --build} $\to$ Full application at \texttt{http://localhost}.
\end{checkpoint}

\section{Exercise 5.4 --- Admin Panel}

\textbf{Tasks:}
\begin{enumerate}
    \item Backend: \texttt{AdminController} with PUT/DELETE endpoints, protected by \texttt{@PreAuthorize}
    \item Frontend: Dashboard with editable product table, delete with confirmation
    \item Route protection with \texttt{adminGuard}
\end{enumerate}

\begin{thinkbox}
Why do we need BOTH \texttt{@PreAuthorize} (backend) AND route guards (frontend)? What if you only protect the frontend?
\end{thinkbox}

% ============================================================
% PART II: SOLUTIONS
% ============================================================
\part{Solutions}

% ============================================================
% CHAPTER 8: SOLUTIONS PHASE 1
% ============================================================
\chapter{Solution --- Phase 1}

\section{Solution 1.1 --- application.yml}

\begin{lstlisting}[style=yamlstyle, title={application.yml}]
server:
  port: 8080

spring:
  datasource:
    url: jdbc:postgresql://localhost:5432/smart_inventory
    username: postgres
    password: secret
    driver-class-name: org.postgresql.Driver
  jpa:
    hibernate:
      ddl-auto: update
    show-sql: true
    properties:
      hibernate:
        format_sql: true
        dialect: org.hibernate.dialect.PostgreSQLDialect
\end{lstlisting}

\section{Solution 1.2 --- Data Model}

\begin{lstlisting}[style=javastyle, title={model/Category.java}]
package com.smartinventory.model;

import jakarta.persistence.*;
import lombok.*;
import java.util.List;

@Entity
@Table(name = "categories")
@Getter @Setter
@NoArgsConstructor @AllArgsConstructor
@Builder
public class Category {
    @Id
    @GeneratedValue(strategy = GenerationType.IDENTITY)
    private Long id;

    @Column(nullable = false, unique = true)
    private String name;

    private String description;

    @OneToMany(mappedBy = "category",
               cascade = CascadeType.ALL,
               fetch = FetchType.LAZY)
    private List<Product> products;
}
\end{lstlisting}

\begin{lstlisting}[style=javastyle, title={model/Product.java}]
package com.smartinventory.model;

import jakarta.persistence.*;
import lombok.*;
import org.hibernate.annotations.*;
import java.math.BigDecimal;
import java.time.LocalDateTime;

@Entity
@Table(name = "products")
@Getter @Setter
@NoArgsConstructor @AllArgsConstructor
@Builder
public class Product {
    @Id
    @GeneratedValue(strategy = GenerationType.IDENTITY)
    private Long id;

    @Column(nullable = false)
    private String name;

    private String description;

    @Column(nullable = false, precision = 10, scale = 2)
    private BigDecimal price;

    @Column(name = "stock_quantity", nullable = false)
    private Integer stockQuantity;

    @Column(name = "image_url")
    private String imageUrl;

    @ManyToOne(fetch = FetchType.LAZY)
    @JoinColumn(name = "category_id", nullable = false)
    private Category category;

    @CreationTimestamp
    @Column(name = "created_at", updatable = false)
    private LocalDateTime createdAt;

    @UpdateTimestamp
    @Column(name = "updated_at")
    private LocalDateTime updatedAt;
}
\end{lstlisting}

\begin{infobox}
\textbf{Why BigDecimal?} \texttt{Double} uses floating-point: \texttt{0.1 + 0.2 = 0.30000000000000004}. For financial calculations, \texttt{BigDecimal} provides exact precision.\\[0.2cm]
\textbf{Why FetchType.LAZY?} With EAGER, loading 100 products loads 100 categories immediately. LAZY loads only when accessed.
\end{infobox}

\section{Solution 1.3 --- Repositories}

\begin{lstlisting}[style=javastyle, title={repository/ProductRepository.java}]
package com.smartinventory.repository;

import com.smartinventory.model.Product;
import org.springframework.data.jpa.repository.JpaRepository;
import java.util.List;

public interface ProductRepository
        extends JpaRepository<Product, Long> {
    List<Product> findByCategoryId(Long categoryId);
}
\end{lstlisting}

\section{Solution 1.4 --- DTOs}

\begin{lstlisting}[style=javastyle, title={dto/response/ProductResponse.java}]
package com.smartinventory.dto.response;

import java.math.BigDecimal;

public record ProductResponse(
    Long id,
    String name,
    String description,
    BigDecimal price,
    Integer stockQuantity,
    String imageUrl,
    String categoryName
) {}
\end{lstlisting}

\begin{lstlisting}[style=javastyle, title={dto/request/ProductRequest.java}]
package com.smartinventory.dto.request;

import java.math.BigDecimal;

public record ProductRequest(
    String name,
    String description,
    BigDecimal price,
    Integer stockQuantity,
    String imageUrl,
    Long categoryId
) {}
\end{lstlisting}

\section{Solution 1.5 --- Service Layer}

\begin{lstlisting}[style=javastyle, title={exception/ResourceNotFoundException.java}]
package com.smartinventory.exception;

public class ResourceNotFoundException
        extends RuntimeException {
    public ResourceNotFoundException(String message) {
        super(message);
    }
}
\end{lstlisting}

\begin{lstlisting}[style=javastyle, title={service/impl/ProductServiceImpl.java}]
package com.smartinventory.service.impl;

import com.smartinventory.dto.request.ProductRequest;
import com.smartinventory.dto.response.ProductResponse;
import com.smartinventory.exception.*;
import com.smartinventory.model.*;
import com.smartinventory.repository.*;
import com.smartinventory.service.ProductService;
import lombok.RequiredArgsConstructor;
import org.springframework.stereotype.Service;
import java.util.List;
import java.util.stream.Collectors;

@Service
@RequiredArgsConstructor
public class ProductServiceImpl
        implements ProductService {

    private final ProductRepository productRepository;
    private final CategoryRepository categoryRepository;

    @Override
    public List<ProductResponse> getAllProducts() {
        return productRepository.findAll()
            .stream()
            .map(this::mapToResponse)
            .collect(Collectors.toList());
    }

    @Override
    public ProductResponse getProductById(Long id) {
        Product product = productRepository.findById(id)
            .orElseThrow(() ->
                new ResourceNotFoundException(
                    "Product not found with id: " + id));
        return mapToResponse(product);
    }

    @Override
    public ProductResponse createProduct(
            ProductRequest request) {
        Category category = categoryRepository
            .findById(request.categoryId())
            .orElseThrow(() ->
                new ResourceNotFoundException(
                    "Category not found with id: "
                    + request.categoryId()));

        Product product = Product.builder()
            .name(request.name())
            .description(request.description())
            .price(request.price())
            .stockQuantity(request.stockQuantity())
            .imageUrl(request.imageUrl())
            .category(category)
            .build();

        Product saved = productRepository.save(product);
        return mapToResponse(saved);
    }

    private ProductResponse mapToResponse(Product p) {
        return new ProductResponse(
            p.getId(), p.getName(),
            p.getDescription(), p.getPrice(),
            p.getStockQuantity(), p.getImageUrl(),
            p.getCategory().getName());
    }
}
\end{lstlisting}

\section{Solution 1.6 --- Controller}

\begin{lstlisting}[style=javastyle, title={controller/ProductController.java}]
package com.smartinventory.controller;

import com.smartinventory.dto.request.ProductRequest;
import com.smartinventory.dto.response.ProductResponse;
import com.smartinventory.service.ProductService;
import lombok.RequiredArgsConstructor;
import org.springframework.http.*;
import org.springframework.web.bind.annotation.*;
import java.util.List;

@RestController
@RequestMapping("/api/v1/products")
@RequiredArgsConstructor
public class ProductController {

    private final ProductService productService;

    @GetMapping
    public ResponseEntity<List<ProductResponse>>
            getAllProducts() {
        return ResponseEntity.ok(
            productService.getAllProducts());
    }

    @GetMapping("/{id}")
    public ResponseEntity<ProductResponse>
            getProductById(@PathVariable Long id) {
        return ResponseEntity.ok(
            productService.getProductById(id));
    }

    @PostMapping
    public ResponseEntity<ProductResponse>
            createProduct(
                @RequestBody ProductRequest request) {
        ProductResponse response =
            productService.createProduct(request);
        return ResponseEntity
            .status(HttpStatus.CREATED)
            .body(response);
    }
}
\end{lstlisting}

\section{Solution 1.7 --- CORS}

\begin{lstlisting}[style=javastyle, title={config/CorsConfig.java}]
package com.smartinventory.config;

import org.springframework.context.annotation.*;
import org.springframework.web.servlet.config.annotation.*;

@Configuration
public class CorsConfig
        implements WebMvcConfigurer {

    @Override
    public void addCorsMappings(CorsRegistry registry) {
        registry.addMapping("/api/**")
            .allowedOrigins("http://localhost:4200")
            .allowedMethods("GET", "POST",
                "PUT", "DELETE", "OPTIONS")
            .allowedHeaders("*")
            .allowCredentials(true)
            .maxAge(3600);
    }
}
\end{lstlisting}

\section{Solution 1.8--1.12 --- Angular Setup}

\begin{lstlisting}[style=tsstyle, title={src/environments/environment.development.ts}]
export const environment = {
  production: false,
  apiUrl: 'http://localhost:8080/api/v1'
};
\end{lstlisting}

\begin{lstlisting}[style=tsstyle, title={src/app/app.config.ts}]
import { ApplicationConfig } from '@angular/core';
import { provideRouter } from '@angular/router';
import { provideHttpClient, withFetch }
  from '@angular/common/http';
import { routes } from './app.routes';

export const appConfig: ApplicationConfig = {
  providers: [
    provideRouter(routes),
    provideHttpClient(withFetch())
  ]
};
\end{lstlisting}

\begin{lstlisting}[style=tsstyle, title={features/products/models/product.model.ts}]
export interface Product {
  id: number;
  name: string;
  description: string;
  price: number;
  stockQuantity: number;
  imageUrl: string;
  categoryName: string;
}

export interface ProductRequest {
  name: string;
  description: string;
  price: number;
  stockQuantity: number;
  imageUrl: string;
  categoryId: number;
}
\end{lstlisting}

\begin{lstlisting}[style=tsstyle, title={features/products/services/product.service.ts}]
import { Injectable } from '@angular/core';
import { HttpClient } from '@angular/common/http';
import { Observable } from 'rxjs';
import { Product, ProductRequest }
  from '../models/product.model';
import { environment }
  from '../../../../environments/environment';

@Injectable({ providedIn: 'root' })
export class ProductService {
  private readonly apiUrl =
    `${environment.apiUrl}/products`;

  constructor(private http: HttpClient) {}

  getAll(): Observable<Product[]> {
    return this.http.get<Product[]>(this.apiUrl);
  }

  getById(id: number): Observable<Product> {
    return this.http.get<Product>(
      `${this.apiUrl}/${id}`);
  }

  create(req: ProductRequest): Observable<Product> {
    return this.http.post<Product>(this.apiUrl, req);
  }
}
\end{lstlisting}

\begin{lstlisting}[style=tsstyle, title={features/products/product-list/product-list.component.ts}]
import { Component, OnInit } from '@angular/core';
import { CommonModule } from '@angular/common';
import { ProductService }
  from '../services/product.service';
import { Product } from '../models/product.model';

@Component({
  selector: 'app-product-list',
  standalone: true,
  imports: [CommonModule],
  templateUrl: './product-list.component.html'
})
export class ProductListComponent implements OnInit {
  products: Product[] = [];
  loading = true;
  error = '';

  constructor(private productService: ProductService) {}

  ngOnInit(): void {
    this.productService.getAll().subscribe({
      next: (data) => {
        this.products = data;
        this.loading = false;
      },
      error: (err) => {
        this.error = 'Failed to load products';
        this.loading = false;
      }
    });
  }
}
\end{lstlisting}

\begin{lstlisting}[style=htmlstyle, title={product-list.component.html}]
<div class="product-list-container">
  <h1>Product Catalog</h1>

  @if (loading) {
    <p>Loading products...</p>
  }

  @if (!loading && !error) {
    <div class="product-grid">
      @for (product of products; track product.id) {
        <div class="product-card">
          <h3>{{ product.name }}</h3>
          <p>{{ product.categoryName }}</p>
          <p>{{ product.price | currency:'USD' }}</p>
          @if (product.stockQuantity === 0) {
            <p class="out-of-stock">Out of Stock</p>
          } @else {
            <p>In Stock: {{ product.stockQuantity }}</p>
          }
        </div>
      } @empty {
        <p>No products found.</p>
      }
    </div>
  }
</div>
\end{lstlisting}

% ============================================================
% CHAPTER 9: SOLUTIONS PHASE 2
% ============================================================
\chapter{Solution --- Phase 2}

\section{Solution 2.1 --- Bean Validation}

\begin{lstlisting}[style=javastyle, title={dto/request/ProductRequest.java (updated)}]
package com.smartinventory.dto.request;

import jakarta.validation.constraints.*;
import java.math.BigDecimal;

public record ProductRequest(
    @NotBlank(message = "Product name is required")
    @Size(max = 100,
          message = "Name must be at most 100 chars")
    String name,

    @Size(max = 500)
    String description,

    @NotNull(message = "Price is required")
    @DecimalMin(value = "0.01",
                message = "Price must be > 0")
    BigDecimal price,

    @NotNull(message = "Stock quantity is required")
    @Min(value = 0,
         message = "Stock cannot be negative")
    Integer stockQuantity,

    String imageUrl,

    @NotNull(message = "Category is required")
    Long categoryId
) {}
\end{lstlisting}

\section{Solution 2.2 --- Global Error Handler}

\begin{lstlisting}[style=javastyle, title={exception/GlobalExceptionHandler.java}]
package com.smartinventory.exception;

import com.smartinventory.dto.response.ErrorResponse;
import org.springframework.http.HttpStatus;
import org.springframework.http.ResponseEntity;
import org.springframework.validation.FieldError;
import org.springframework.web.bind.MethodArgumentNotValidException;
import org.springframework.web.bind.annotation.*;
import java.time.LocalDateTime;
import java.util.*;

@RestControllerAdvice
public class GlobalExceptionHandler {

    @ExceptionHandler(
        MethodArgumentNotValidException.class)
    public ResponseEntity<ErrorResponse>
            handleValidation(
                MethodArgumentNotValidException ex) {
        Map<String, String> errors = new HashMap<>();
        ex.getBindingResult().getFieldErrors()
            .forEach((FieldError fe) ->
                errors.put(fe.getField(),
                           fe.getDefaultMessage()));

        ErrorResponse response = new ErrorResponse(
            LocalDateTime.now(),
            HttpStatus.BAD_REQUEST.value(),
            "Validation Failed", errors);

        return ResponseEntity
            .status(HttpStatus.BAD_REQUEST)
            .body(response);
    }

    @ExceptionHandler(
        ResourceNotFoundException.class)
    public ResponseEntity<ErrorResponse>
            handleNotFound(
                ResourceNotFoundException ex) {
        ErrorResponse response = new ErrorResponse(
            LocalDateTime.now(),
            HttpStatus.NOT_FOUND.value(),
            ex.getMessage(), null);
        return ResponseEntity
            .status(HttpStatus.NOT_FOUND)
            .body(response);
    }

    @ExceptionHandler(Exception.class)
    public ResponseEntity<ErrorResponse>
            handleGeneric(Exception ex) {
        ex.printStackTrace();
        ErrorResponse response = new ErrorResponse(
            LocalDateTime.now(),
            HttpStatus.INTERNAL_SERVER_ERROR.value(),
            "An unexpected error occurred", null);
        return ResponseEntity
            .status(HttpStatus.INTERNAL_SERVER_ERROR)
            .body(response);
    }
}
\end{lstlisting}

\section{Solution 2.3--2.5 --- Angular Forms}

\begin{lstlisting}[style=tsstyle, title={product-form.component.ts}]
import { Component } from '@angular/core';
import { CommonModule } from '@angular/common';
import { ReactiveFormsModule, FormGroup,
  FormControl, Validators } from '@angular/forms';
import { Router } from '@angular/router';
import { ProductService }
  from '../services/product.service';
import { HttpErrorResponse }
  from '@angular/common/http';
import { priceFormatValidator }
  from '../validators/price.validator';

@Component({
  selector: 'app-product-form',
  standalone: true,
  imports: [CommonModule, ReactiveFormsModule],
  templateUrl: './product-form.component.html'
})
export class ProductFormComponent {
  serverErrors: { [key: string]: string } = {};
  submitting = false;

  form = new FormGroup({
    name: new FormControl('', [
      Validators.required,
      Validators.minLength(2),
      Validators.maxLength(100)
    ]),
    description: new FormControl('',
      [Validators.maxLength(500)]),
    price: new FormControl<number | null>(null, [
      Validators.required,
      Validators.min(0.01),
      priceFormatValidator
    ]),
    stockQuantity: new FormControl<number | null>(
      null, [Validators.required, Validators.min(0)]),
    imageUrl: new FormControl(''),
    categoryId: new FormControl<number | null>(
      null, [Validators.required])
  });

  constructor(
    private productService: ProductService,
    private router: Router
  ) {
    this.form.valueChanges.subscribe(() => {
      this.serverErrors = {};
    });
  }

  onSubmit(): void {
    this.form.markAllAsTouched();
    if (this.form.invalid) return;

    this.submitting = true;
    this.productService.create(
      this.form.value as any
    ).subscribe({
      next: () =>
        this.router.navigate(['/products']),
      error: (err: HttpErrorResponse) => {
        this.submitting = false;
        if (err.status === 400 && err.error?.errors) {
          this.serverErrors = err.error.errors;
        }
      }
    });
  }

  hasError(field: string, type: string): boolean {
    const ctrl = this.form.get(field);
    return !!(ctrl?.hasError(type) && ctrl?.touched);
  }
}
\end{lstlisting}

\begin{lstlisting}[style=tsstyle, title={validators/price.validator.ts}]
import { AbstractControl, ValidationErrors }
  from '@angular/forms';

export function priceFormatValidator(
    control: AbstractControl
): ValidationErrors | null {
  const value = control.value;
  if (!value) return null;

  const str = value.toString();
  const dotIndex = str.indexOf('.');
  if (dotIndex !== -1) {
    const decimals = str.length - dotIndex - 1;
    if (decimals > 2) return { invalidPrice: true };
  }
  return null;
}
\end{lstlisting}

% ============================================================
% CHAPTER 10: SOLUTIONS PHASE 3
% ============================================================
\chapter{Solution --- Phase 3}

\section{Solution 3.1--3.2 --- Pagination \& Specifications}

\begin{lstlisting}[style=javastyle, title={repository/ProductSpecification.java}]
package com.smartinventory.repository;

import com.smartinventory.model.Product;
import org.springframework.data.jpa.domain.Specification;
import java.math.BigDecimal;

public class ProductSpecification {

    public static Specification<Product>
            hasName(String name) {
        return (root, query, cb) -> {
            if (name == null || name.isBlank())
                return null;
            return cb.like(
                cb.lower(root.get("name")),
                "%" + name.toLowerCase() + "%");
        };
    }

    public static Specification<Product>
            hasCategory(Long categoryId) {
        return (root, query, cb) -> {
            if (categoryId == null) return null;
            return cb.equal(
                root.get("category").get("id"),
                categoryId);
        };
    }

    public static Specification<Product>
            hasPriceRange(
                BigDecimal min, BigDecimal max) {
        return (root, query, cb) -> {
            if (min == null && max == null)
                return null;
            if (min != null && max != null)
                return cb.between(
                    root.get("price"), min, max);
            if (min != null)
                return cb.greaterThanOrEqualTo(
                    root.get("price"), min);
            return cb.lessThanOrEqualTo(
                root.get("price"), max);
        };
    }
}
\end{lstlisting}

\section{Solution 3.3--3.4 --- Signals \& RxJS}

\begin{lstlisting}[style=tsstyle, title={product-list.component.ts (with Signals \& RxJS)}]
import { Component, OnInit, OnDestroy,
  signal, computed } from '@angular/core';
import { CommonModule } from '@angular/common';
import { ReactiveFormsModule, FormControl }
  from '@angular/forms';
import { Subject, takeUntil } from 'rxjs';
import { debounceTime, distinctUntilChanged,
  switchMap, tap } from 'rxjs/operators';
import { ProductService }
  from '../services/product.service';
import { Product, PageResponse }
  from '../models/product.model';

@Component({
  selector: 'app-product-list',
  standalone: true,
  imports: [CommonModule, ReactiveFormsModule],
  templateUrl: './product-list.component.html'
})
export class ProductListComponent
    implements OnInit, OnDestroy {

  searchControl = new FormControl('');

  products = signal<Product[]>([]);
  loading = signal(true);
  currentPage = signal(0);
  totalPages = signal(0);
  pageSize = signal(10);

  isEmpty = computed(() =>
    this.products().length === 0
    && !this.loading());

  private destroy$ = new Subject<void>();

  constructor(
    private productService: ProductService) {}

  ngOnInit(): void {
    this.loadProducts('');

    this.searchControl.valueChanges.pipe(
      debounceTime(300),
      distinctUntilChanged(),
      tap(() => {
        this.loading.set(true);
        this.currentPage.set(0);
      }),
      switchMap(term =>
        this.productService.search(
          term || '', 0, this.pageSize())),
      takeUntil(this.destroy$)
    ).subscribe({
      next: (res: PageResponse<Product>) => {
        this.products.set(res.content);
        this.totalPages.set(res.totalPages);
        this.loading.set(false);
      },
      error: () => this.loading.set(false)
    });
  }

  loadProducts(search: string): void {
    this.loading.set(true);
    this.productService.search(
      search, this.currentPage(), this.pageSize()
    ).pipe(takeUntil(this.destroy$))
    .subscribe({
      next: (res) => {
        this.products.set(res.content);
        this.totalPages.set(res.totalPages);
        this.loading.set(false);
      },
      error: () => this.loading.set(false)
    });
  }

  ngOnDestroy(): void {
    this.destroy$.next();
    this.destroy$.complete();
  }
}
\end{lstlisting}

% ============================================================
% CHAPTER 11: SOLUTIONS PHASE 4
% ============================================================
\chapter{Solution --- Phase 4}

\section{Solution 4.1 --- User Entity}

\begin{lstlisting}[style=javastyle, title={model/User.java (implements UserDetails)}]
package com.smartinventory.model;

import jakarta.persistence.*;
import lombok.*;
import org.hibernate.annotations.CreationTimestamp;
import org.springframework.security.core.*;
import org.springframework.security.core.authority.*;
import org.springframework.security.core.userdetails.*;
import java.time.LocalDateTime;
import java.util.*;

@Entity @Table(name = "users")
@Getter @Setter @NoArgsConstructor
@AllArgsConstructor @Builder
public class User implements UserDetails {
    @Id @GeneratedValue(strategy =
        GenerationType.IDENTITY)
    private Long id;

    @Column(name = "first_name", nullable = false)
    private String firstName;

    @Column(name = "last_name", nullable = false)
    private String lastName;

    @Column(nullable = false, unique = true)
    private String email;

    @Column(nullable = false)
    private String password;

    @Enumerated(EnumType.STRING)
    @Column(nullable = false)
    private Role role;

    @CreationTimestamp
    @Column(name = "created_at", updatable = false)
    private LocalDateTime createdAt;

    @Override
    public Collection<? extends GrantedAuthority>
            getAuthorities() {
        return List.of(new SimpleGrantedAuthority(
            "ROLE_" + role.name()));
    }

    @Override
    public String getUsername() { return email; }
    @Override
    public boolean isAccountNonExpired() {
        return true; }
    @Override
    public boolean isAccountNonLocked() {
        return true; }
    @Override
    public boolean isCredentialsNonExpired() {
        return true; }
    @Override
    public boolean isEnabled() { return true; }
}
\end{lstlisting}

\section{Solution 4.2--4.3 --- JWT \& Security Config}

See the complete solutions in the Markdown version of this TP. The key patterns are:

\begin{enumerate}
    \item \texttt{JwtService}: Uses \texttt{io.jsonwebtoken} to generate/validate tokens
    \item \texttt{JwtAuthenticationFilter}: Extends \texttt{OncePerRequestFilter}, extracts Bearer token, validates, sets \texttt{SecurityContextHolder}
    \item \texttt{SecurityConfig}: Disables CSRF, sets STATELESS sessions, permits \texttt{/auth/**}, requires ADMIN for \texttt{/admin/**}
\end{enumerate}

\section{Solution 4.5--4.7 --- Angular Auth}

\begin{lstlisting}[style=tsstyle, title={core/interceptors/auth.interceptor.ts}]
import { HttpInterceptorFn }
  from '@angular/common/http';
import { inject } from '@angular/core';
import { StorageService }
  from '../services/storage.service';

export const authInterceptor: HttpInterceptorFn =
    (req, next) => {
  if (req.url.includes('/auth/'))
    return next(req);

  const storage = inject(StorageService);
  const token = storage.getToken();

  if (token) {
    const cloned = req.clone({
      setHeaders: {
        Authorization: `Bearer ${token}`
      }
    });
    return next(cloned);
  }

  return next(req);
};
\end{lstlisting}

\begin{lstlisting}[style=tsstyle, title={core/guards/auth.guard.ts}]
import { CanActivateFn, Router }
  from '@angular/router';
import { inject } from '@angular/core';
import { AuthService }
  from '../services/auth.service';

export const authGuard: CanActivateFn = () => {
  const auth = inject(AuthService);
  const router = inject(Router);

  if (auth.isLoggedIn()) return true;
  return router.createUrlTree(['/auth/login']);
};
\end{lstlisting}

% ============================================================
% CHAPTER 12: SOLUTIONS PHASE 5
% ============================================================
\chapter{Solution --- Phase 5}

\section{Solution 5.1 --- N+1 Fix}

\begin{lstlisting}[style=javastyle, title={ProductRepository with @EntityGraph}]
@EntityGraph(attributePaths = {"category"})
Page<Product> findAll(
    Specification<Product> spec, Pageable pageable);

@EntityGraph(attributePaths = {"category"})
Page<Product> findAll(Pageable pageable);
\end{lstlisting}

\begin{infobox}
\textbf{Before:} 1 query for products + N queries for categories = N+1 queries\\
\textbf{After:} 1 query with LEFT JOIN = 1 query total
\end{infobox}

\section{Solution 5.3 --- Docker}

\begin{lstlisting}[style=dockerstyle, title={docker-compose.yml}]
version: '3.8'

services:
  db:
    image: postgres:16-alpine
    environment:
      POSTGRES_DB: smart_inventory
      POSTGRES_USER: postgres
      POSTGRES_PASSWORD: ${DB_PASSWORD}
    volumes:
      - postgres_data:/var/lib/postgresql/data
    healthcheck:
      test: ["CMD-SHELL",
             "pg_isready -U postgres"]
      interval: 10s
      timeout: 5s
      retries: 5

  api:
    build: ./smart-inventory-api
    environment:
      SPRING_DATASOURCE_URL:
        jdbc:postgresql://db:5432/smart_inventory
      SPRING_DATASOURCE_USERNAME: postgres
      SPRING_DATASOURCE_PASSWORD: ${DB_PASSWORD}
    ports:
      - "8080:8080"
    depends_on:
      db:
        condition: service_healthy

  ui:
    build: ./smart-inventory-ui
    ports:
      - "80:80"
    depends_on:
      - api

volumes:
  postgres_data:
\end{lstlisting}

% ============================================================
% WEEKLY SCHEDULE
% ============================================================
\chapter{Weekly Time Distribution}

\section{Daily Schedule (2 hours/day)}

\begin{longtable}{clp{9cm}}
\toprule
\textbf{Day} & \textbf{Focus} & \textbf{Activities} \\
\midrule
\endhead
1 & \faIcon{book} + \faIcon{server} & 30min: Read docs/video. 90min: Implement API endpoint \\[0.3cm]
2 & \faIcon{server} & 30min: Read docs. 90min: Complete backend, validation, errors \\[0.3cm]
3 & \faIcon{book} + \faIcon{laptop} & 30min: Read Angular docs. 90min: Build service \& component \\[0.3cm]
4 & \faIcon{laptop} & 30min: Read docs. 90min: Complete UI, connect to backend \\[0.3cm]
5 & \faIcon{search} & Full 2h: Refactoring, debugging (Network Tab, SQL logs) \\[0.3cm]
6 & \faIcon{flask} & Full 2h: Add a feature using only official docs \& error messages \\
\bottomrule
\end{longtable}

% ============================================================
% RESOURCES
% ============================================================
\chapter{Essential Resources}

\begin{longtable}{p{4cm}l}
\toprule
\textbf{Resource} & \textbf{Link} \\
\midrule
\endhead
Spring Boot Reference & \url{https://docs.spring.io/spring-boot/reference/} \\[0.3cm]
Spring Data JPA & \url{https://docs.spring.io/spring-data/jpa/reference/} \\[0.3cm]
Angular Documentation & \url{https://angular.dev} \\[0.3cm]
RxJS Documentation & \url{https://rxjs.dev} \\[0.3cm]
JWT Introduction & \url{https://jwt.io/introduction} \\[0.3cm]
Docker Compose & \url{https://docs.docker.com/compose/} \\
\bottomrule
\end{longtable}

\vfill

\begin{center}
\rule{0.5\textwidth}{0.5pt}\\[0.5cm]
\textit{End of TP --- Smart Inventory \& Order Platform}\\[0.2cm]
\textcopyright\ 2026 Yassine Elkhamlichi
\end{center}

\end{document}
